\section{詳細解説}

\subsection{腕 1 本の振りで得られる回転の地図}

手足が 1 本なら、胴体の回転は関節の始点と終点だけで決まる（$\theta(q_e)-\theta(q_s)=-\int_{q_s}^{q_e}\Gamma\,dq$）。
図~\ref{fig:D2_swing_map.png} は、始点と終点のすべての組み合わせについて胴体の回転を示したものである。
最大の首振りは約 32°（腕を下側から上側へ大きく回す場合）で、Sim01 の「下 → 上の 180° で $-16.0^\circ$」はその一部である。
\textbf{腕 1 本では、どんな振り方をしてもこの地図の範囲を超えられない}。

\DFIG{D2_swing_map.png}{腕 1 本の振りによる胴体の回転 [°]。横軸が始点、縦軸が終点。経路によらず、始点と終点だけで決まる。}

\subsection{質量比・長さ比のスイープの全結果}

Sim01 で腕の質量比 $\mu$（腕の質量 / 全質量）と長さ比（腕の長さ / 胴体の高さ）を変えたときの、180° の振りによる回転角を示す。
回転角は質量比と長さ比のどちらにもほぼ比例して増え、人間の腕（質量比 0.05〜0.06、長さ比 0.7）では十数度である。

\DTAB{sim01_sweep}

\subsection{曲率マップと手足の質量}

腕と脚の質量を変えて、曲率マップと 1 周あたりの回転の上界（関節可動域内で同じ符号の曲率をすべて囲んだ場合）を比べた（図~\ref{fig:D3_curvature_mass.png}、総質量 76 kg）。

\begin{center}\begin{tabular}{lll}
\toprule
腕 & 脚 & 1 周あたりの上界 \\ \midrule
4 kg & 6 kg & 6.41° \\
4 kg & 12 kg（人間型） & 9.60° \\
4 kg & 24 kg & 12.91° \\
8 kg & 12 kg & 16.90° \\
\bottomrule
\end{tabular}\end{center}

腕を 2 倍にする方が、脚を 2 倍にするより効果が大きい。腕は肩の位置が胴体重心から遠く、可動域も広い（$-150^\circ$〜$150^\circ$）ためである。
曲率が負の領域（青）は、どの場合も腕を前方〜やや下、脚を前方に上げた付近に広がり、質量を増やすと領域の形はあまり変わらずに濃く（曲率が大きく）なる。

\DFIG{D3_curvature_mass.png}{手足の質量と曲率マップ。破線は $F=0$。タイトルの値は関節可動域内の 1 周あたりの回転の上界。}

\subsection{ストークスの定理による検証の詳細}

Sim02 では、腕を $[-90^\circ, 90^\circ]$、脚を $[-90^\circ, -30^\circ]$ の長方形ループで回した。
\begin{itemize}
\item \textbf{数値積分}：各辺を最小躍度のプロファイルで 1 s ずつ動かし、$\dot\theta=-\Gamma\cdot\dot q$ を適応刻みで積分した（相対誤差 $10^{-10}$）。
\item \textbf{面積分}：長方形の内部で曲率 $F$（中心差分）を 2 重積分した。
\end{itemize}
両者は $+2.9778^\circ$ で一致した。数値積分の結果は各辺の速さの配分によらないことも確認でき、時間に依存しないこと（第~\ref{ch:theory}章）の数値的な確認にもなっている。

\section*{考察課題}
\begin{exercise}
図~\ref{fig:D3_curvature_mass.png} の曲率マップで、Sim02 の長方形ループは曲率の弱い領域に置かれていた。同じ大きさの長方形をどこに置けば 1 周あたりの回転が最大になるか、地図から推定せよ。
\end{exercise}
\begin{exercise}
腕の可動域を $[-90^\circ, 90^\circ]$ に狭めると、1 周あたりの上界はどう変わるか。曲率マップのどの部分が使えなくなるかから考えよ。
\end{exercise}
